\section{Data and Experimental Setup}

\label{sec:data}
% ══════════════════════════════════════════════════════════════════════════════

\subsection{Species Occurrence Data}

We compiled occurrence records from multiple sources:

\begin{itemize}
  \item \textbf{GBIF} \citep{gbif2024}: 12.4 million georeferenced
    records for 847 terrestrial vertebrate species.
  \item \textbf{eBird} \citep{sullivan2009ebird}: 8.2 million
    structured bird survey records.
  \item \textbf{Camera trap networks}: 1.3 million detection records
    from TEAM \citep{jansen2014team} and Wildlife Insights
    \citep{ahumada2020wildlife}.
\end{itemize}

After spatial thinning (minimum 1\,km separation) and temporal filtering
(2000--2023), the final dataset contains 4.7 million occurrence records.

\subsection{Environmental Variables}

Bioclimatic variables were sourced from WorldClim v2.1
\citep{fick2017worldclim} at 30 arc-second resolution. Topographic
variables were derived from SRTM v4.1 \citep{jarvis2008srtm}. Land cover
was from ESA WorldCover 2021 \citep{zanaga2022esa}. MODIS vegetation
indices were averaged over 2018--2022.

\subsection{Climate Projections}

We used downscaled CMIP6 projections from WorldClim
\citep{fick2017worldclim} for four SSP scenarios at three time periods
(2041--2060, 2061--2080, 2081--2100). Twelve GCMs were selected based
on performance evaluation by \citet{arias2021climate}.

\subsection{Study Regions}

We evaluate across six biogeographic realms:
\begin{enumerate}
  \item Neotropical (Amazon basin focus)
  \item Afrotropical (East Africa highlands)
  \item Indo-Malayan (Southeast Asian tropics)
  \item Palearctic (European temperate forests)
  \item Nearctic (North American boreal-temperate transition)
  \item Australasian (Eastern Australia)
\end{enumerate}

\subsection{Baseline Models}

\begin{itemize}
  \item \textbf{MaxEnt} \citep{phillips2006maximum}: Standard
    configuration with linear, quadratic, hinge, and threshold features.
  \item \textbf{Random Forest (RF)}: 500 trees, following
    \citet{cutler2007random}.
  \item \textbf{BRT}: Boosted regression trees \citep{elith2008working}
    with 5000 trees and interaction depth 5.
  \item \textbf{DeepSDM}: 4-layer MLP with 256 hidden units per layer
    \citep{botella2018deep}.
  \item \textbf{CNN-SDM}: Convolutional model on environmental raster
    patches \citep{deneu2021convolutional}.
\end{itemize}

\subsection{Evaluation Protocol}

We employ two validation strategies:
\begin{enumerate}
  \item \textbf{Spatial block cross-validation}: Study area divided into
    $10 \times 10$ blocks, with blocks assigned to 5 folds ensuring
    minimum 200\,km separation between train and test
    \citep{roberts2017cross}.
  \item \textbf{Temporal validation}: Models trained on 2000--2015 data,
    validated on 2016--2023 independent resurvey data from standardized
    monitoring programs.
\end{enumerate}

Metrics include AUC, True Skill Statistic (TSS), and Boyce index for
continuous suitability evaluation \citep{hirzel2006evaluating}.


% ══════════════════════════════════════════════════════════════════════════════
